\chapter{Introducción}\label{intro}

\textit{\textbf{El presente documento sirve de plantilla para la presentación del informe final del trabajo de grado de la Maestría en Ingeniería del Internet de las Cosas. El formato no puede ser modificado (márgenes, numeración, tamaño de fuentes, referencias, etc.) y la estructura de los capítulos también deberá mantenerse.}}

\textit{\textbf{La extensión máxima nunca deberá exceder las 60 páginas en total, incluyendo todas las diferentes secciones, portadas y referencias. Los anexos son opcionales, no se consideran parte central del documento y no se contarán dentro de la extensión máxima.}}\\[1.0cm]

En la Introducción se presenta la problematica que se desarrolló durante el trabajo de grado, así como la justificación y la motivación. Así mismo, se presentan los objetivos, alcances, limitaciones y la metodología utilizada. Se deberá hacer un análisis general de como el trabajo desarrollado aporta en el avance del campo de aplicación específico. No se debe confundir con el resumen.

\textcolor{red}{Se recomienda tener una extensión de $3$ páginas.}


